\section{Problem, notation, and scope}
\label{sec:problem-and-scope}

Write \(\T=\mathbb R/\mathbb Z\), with normalized Haar measure \(\mu\).  For an
integer \(m\geq 3\), a measurable set \(A\subset\T\) is \emph{\(m\)-sum-free}
when the equation
\[
    x+y=mz
\]
has no solution \((x,y,z)\in A^3\).  Equivalently,
\[
    (A+A)\cap mA=\varnothing
    \qquad\text{or}\qquad
    0\notin A+A-mA.
\]
For a compact abelian group \(G\), set
\[
    d_m(G)=\sup\{\mu_G(A):A\subset G\text{ is measurable and }m\text{-sum-free}\}.
\]
The main objective is to determine, or sharply bound, \(d_m(\T)\), with
particular emphasis on \(m=3\) and on the behavior as \(m\to\infty\).

Candela and Sisask prove that, for every fixed finite family of linear forms in
at least three variables, the corresponding extremal densities in \(\Zp\)
converge to the density on \(\T\) as \(p\to\infty\) through primes
\cite[Theorem~1.3]{CandelaSisask2011}.  Applied to \(x+y-mz\), this gives
\[
    \lim_{\substack{p\to\infty\\p\ \mathrm{prime}}}d_m(\Zp)=d_m(\T).
\]
Consequently, a circle construction gives asymptotic finite-group lower
bounds, and uniform finite-group upper bounds can be transferred back to the
circle.  The project therefore moves between \(\T\), where unions of intervals
are geometric, and \(\Zp\), where inverse additive theorems are available.

\subsection{Status at a glance}

The word \emph{computed} below means a numerical MILP result at the stated
solver tolerance.  It does not mean an exact theorem unless a separate proof
is supplied.

\begin{center}
\small
\begin{tabularx}{\textwidth}{@{}>{\raggedright\arraybackslash}p{0.20\textwidth}
>{\raggedright\arraybackslash}p{0.16\textwidth}
>{\raggedright\arraybackslash}X@{}}
\toprule
Topic & Status & Current conclusion \\
\midrule
\(m=3\), two closed intervals & \statusproved & Measure at least \(1/5\)
forces \(A+A-3A=\T\); see Section~\ref{sec:closed-two-interval-proof}. \\
\(m=3\), two open intervals at measure \(1/5\) & \statusproved & The only
non-covering endpoint configuration, up to the stated symmetries, is
\((0,1/10)\cup(1/2,3/5)\); see Sections~\ref{sec:open-two-interval-r5}
and~\ref{sec:open-two-interval-r6}. \\
\(m=3\), \(3\leq N\leq 8\) intervals & \statuscomputed & Reported MILP runs
give optimum \(1/5\).  Reproducible solver artifacts for all these runs still
need to be archived. \\
Omitted interval interactions & \statuscomputed & For \(N=2,3,4,5\), reported
relaxed MILPs identify maximal families of terms \(I_i+I_j-3I_\ell\) that can
be omitted without changing the numerical \(1/5\) optimum. \\
\(m=4\) & \statusproved & Three explicit open intervals give density
\(11/64>1/6\); see Section~\ref{sec:explicit-m4-m5}. \\
\(m=5\) & \statusproved & Two explicit open intervals give density
\(6/35>1/7\); see Section~\ref{sec:explicit-m4-m5}. \\
General \(m\) lower bound & \statusproved & Equally spaced intervals give
\(\lfloor(m+2)^2/8\rfloor/(m(m+2))>1/8\); see
Section~\ref{sec:equally-spaced-construction}. \\
General \(m\) upper bound & literature & For each fixed \(m\geq3\),
\(d_m(\T)\leq1/3.1955\), via the finite-group theorem and transference
\cite[Theorem~1.5]{CandelaGonzalezSanchezGrynkiewicz2021}. \\
Exact asymptotics as \(m\to\infty\) & \statusopen & The present lower bound
tends to \(1/8\), while the cited upper bound is uniform and stays above
\(0.31\). \\
\bottomrule
\end{tabularx}
\end{center}

\subsection{How to read this document}

Sections containing complete proofs are intended to be stable.  Sections on
the MILP, excluded interactions, and tiny added intervals describe a discovery
tool and must retain solver metadata before their output can serve as a
computer-assisted certificate.  The final roadmap states the missing steps
explicitly so that experimental observations are not silently promoted to
claims.
